% TOP_PROBLEM: {"id": 1500010, "title": "Algebraic Stories — Hilbert series of other classes of ideals", "category_id": 4, "category": {"id": 4, "name": "algebra", "display_name": "Algebra", "description": "Group theory, ring theory, field theory, and algebraic structures.", "slug": "algebra", "order_index": 4, "created_at": "2026-07-31T15:26:25.671Z"}, "status": "open"}
% FIRST_POSTED: null
% ENTRY_KIND: statement_only
% Imported from: batch06.tex; UnsolvedMath ID 1500010
% Compile twice: lualatex 1500010.tex
\documentclass[11pt]{article}
\usepackage{fontspec}
\setmainfont{Latin Modern Roman}
\usepackage{amsmath,amssymb,mathtools,unicode-math}
\setmathfont{Latin Modern Math}
\usepackage{xurl,hyperref}
\hypersetup{colorlinks=true,urlcolor=blue,pdftitle={UnsolvedMath Problem 1500010}}
\newcommand{\sref}[2]{\hyperlink{source:#1:#2}{[S#2]}}
\newcommand{\cataloguescope}[1]{\paragraph{Mathematical question.}#1}
\begin{document}
% UnsolvedMath ID 1500010; source catalogue code AMR-014-0010
\setcounter{section}{1500009}
\section{Algebraic Stories — Hilbert series of other classes of ideals}\label{1500010}
{\small\sffamily UnsolvedMath ID 1500010 · Algebra}
\subsection{Definitions and mathematical statement}

\paragraph{Shared notation.}$\mathbb N=\{1,2,\ldots\}$, and $\mathbb Z,\mathbb Q,\mathbb R,\mathbb C$ denote the integers, rationals, reals and complex numbers. Any other conventions are specified locally.


Let $S=\mathbb C[x_1,\ldots,x_n]$ with its usual total-degree grading. For a homogeneous quotient $R$, define
\[\operatorname{Hilb}_R(t)=\sum_{e\geq0}(\dim_{\mathbb C}R_e)t^e.
\]
A generic tuple of forms has its coefficient vector in a suitable nonempty Zariski open subset of the product of the prescribed homogeneous pieces.For a series $A(t)=\sum_{e\geq0}a_et^e$, use $[A(t)]_+$ for its truncation at the first non-positive coefficient: if $e_0=\min\{e:a_e\leq0\}$, retain the terms with $e<e_0$ and replace all later terms by zero; if no such $e_0$ exists, retain the entire series.Let $d,n,r\geq1$ and let $\mu=(\mu_1,\ldots,\mu_q)$ be a partition of $d$ into positive parts. For $i=1,\ldots,r$ choose generic linear forms $\ell_{i,1},\ldots,\ell_{i,q}$, distinct within each product, and set
\[I=\left(\prod_{j=1}^q\ell_{1,j}^{\mu_j},\ldots,\prod_{j=1}^q\ell_{r,j}^{\mu_j}\right).
\]
This is a generic $\mu$-power ideal.
\paragraph{Statement recorded in the supplied source.}
For every $\mu\neq(d)$, is the Hilbert function of $S/I$ the coefficient sequence of
\[\left[\frac{(1-t^d)^r}{(1-t)^n}\right]_+?\] \sref{1500010}{2}
\subsection{Short English statement}
Do products of powers of distinct generic linear forms have the predicted generic-forms Hilbert function?
\subsection{Sources}
\begin{description}
\item[\hypertarget{source:1500010:1}{S1}] ULAM.AI and the UnsolvedMath Contributors, UnsolvedMath: A Curated Collection of Open Mathematics Problems, record 1500010 (AMR-014-0010). CC BY 4.0. \url{https://huggingface.co/datasets/ulamai/UnsolvedMath}.
\item[\hypertarget{source:1500010:2}{S2}] R. Fröberg, S. Lundqvist, A. Oneto and B. Shapiro, \emph{Algebraic Stories from One and from the Other Pockets}, arXiv:1801.01692 (2018), §2.2, Problem F. \url{https://arxiv.org/abs/1801.01692}.
\end{description}
\label{end:1500010}
\end{document}
